\documentclass[12pt]{article}

\usepackage[margin=1in]{geometry}
\usepackage{booktabs}
\usepackage{amsmath}
\usepackage{setspace}
\onehalfspacing

\title{Fixed Effects Illustration:\\Minimum Wage and Unemployment}
\author{ECON 4999B Course Material}
\date{\today}

\begin{document}

\maketitle

\section{Data}

We use a balanced panel of 26 regions (A--Z) and 27 years (2000--2026), for
$26 \times 27 = 702$ observations. The outcome variable is
\texttt{unemployment\_rate} (percent) and the regressor of interest is
\texttt{min\_wage} (dollars).

\section{Specifications}

We estimate
\begin{equation}
    \text{unemployment\_rate}_{it} = \alpha + \beta\,\text{min\_wage}_{it}
    + \mu_i + \lambda_t + \varepsilon_{it},
\end{equation}
where $\mu_i$ denotes region fixed effects and $\lambda_t$ denotes year fixed
effects. The four columns differ by which fixed effects are included:
\begin{enumerate}
    \item[(1)] pooled OLS (no fixed effects);
    \item[(2)] year fixed effects only;
    \item[(3)] region fixed effects only;
    \item[(4)] region and year fixed effects (two-way FE).
\end{enumerate}
All models use heteroskedasticity-robust standard errors (HC1).

\section{Results}

\begin{table}[htbp]
\centering
\caption{Effect of minimum wage on unemployment under alternative specifications}
\label{tab:fe_results}
\begin{tabular}{lcccc}
\toprule
 & (1) & (2) & (3) & (4) \\
\midrule
\multicolumn{5}{l}{\textit{Dependent variable: unemployment rate (\%)}} \\
\addlinespace
Min wage & -0.371*** & -0.530*** & 0.039 & 0.174** \\
 & (0.023) & (0.017) & (0.030) & (0.081) \\
\addlinespace
Fixed effects & None & Year & Region & Region, Year \\
Observations & \multicolumn{4}{c}{702} \\
$R^2$ & 0.237 & 0.702 & 0.588 & 0.941 \\
Adjusted $R^2$ & 0.236 & 0.690 & 0.572 & 0.936 \\
\bottomrule
\end{tabular}

\medskip
\footnotesize
\textit{Notes:} HC1 robust standard errors in parentheses.
$^{*}p<0.10$, $^{**}p<0.05$, $^{***}p<0.01$.
\end{table}

\section{Discussion}

Without fixed effects, the estimated minimum-wage coefficient reflects both
within-region variation and persistent cross-region differences. Regions with
historically higher minimum wages also differ in other ways (industry mix,
demographics, labor-market institutions), so the pooled estimate can be
misleading.

Adding region fixed effects (column 3) uses only within-region changes in the
minimum wage, removing time-invariant regional confounders. Adding year fixed
effects (column 2) removes common macro shocks shared by all regions in a given
year. Two-way fixed effects (column 4) combine both: identification comes from
region-specific deviations from national trends.

In this synthetic dataset, the pooled estimate is negative because high-wage
regions have lower baseline unemployment for reasons unrelated to minimum-wage
changes. Once region fixed effects are included, the coefficient turns positive,
consistent with the within-region relationship built into the data-generating
process.

\end{document}
